\documentclass[main.tex]{subfiles}       

\begin{document}

[Ahuj68] Ahuja, D.V., 和 Coons, S.A., 建筑和展示的几何，IBM 系统杂志，第 3-4 期，188-205 页，1968 年。

[Akim70] Akima, H., 基于局部程序的新插值和光滑曲线拟合方法，ACM 杂志，第 17 卷，589-602 页，1970 年。

[Akma92] Akman, V., 和 Arslan, A., 在拓扑图画书中使用所有图形元素进行扫掠，计算机与图形，第 16 卷，第 3 期，273-281 页，1992 年。

[Ball74] Ball, A., CONSURF. 第 1 部分:锥形放样平铺的介绍，CAD，第 6 卷，第 4 期，243-249 页，1974 年。

[Ball75] Ball, A., CONSURF. 第 2 部分:算法描述，CAD，第 7 卷，第 4 期，237-242 页，1975 年。

[Ball77] Ball, A., CONSURF. 第3部分:程序的使用，CAD，第9卷，第1期，页码9-12，1977年。

[Baly94] Ballay, H. 和 Storn, R., 一种检查C语言编码规范的工具，C/C++用户期刊，第12卷，第27期，页码41-50，1994年。

[Barn93] Barnhill, R., Coons'补丁和凸组合，见[Pieg93]。

[Barr83] Barr, A.H., 固体原始体的全局和局部变形，SIGGRAPH 83教程笔记，底特律，1983年。

[Bars89] Barsky, B.A. 和 DeRose, T.D., 参数曲线的几何连续性:三种等价表征，IEEE计算机图形与应用，第9卷，第6期，页码60-68，1989年。

[Bars90] Barsky, B.A. 和 DeRose, T.D., 参数曲线的几何连续性:几何连续样条的构造，IEEE计算机图形与应用，第10卷，第1期，页码60-68，1990年。

[Bart87] Bartels, R.H., Beatty, J.C. 和 Barsky, B.A., 计算机图形学和几何建模中样条的介绍，加利福尼亚州圣马特奥:摩根·考夫曼，1987年。

[Beac91] Beach, R.C., 计算机辅助设计的曲线和曲面的介绍，纽约:范诺斯特兰德·雷恩霍尔德，1991年。

[Bern12] Bernstein, S.N.，基于概率计算的维尔斯特拉斯定理证明，Commun. Soc. Math. Khrakow，第12卷，第2期，页码 \(1 - 2,{1912}\) 。

[Bezi72] Bézier, P.E.，《数值控制:数学与应用》，纽约:约翰·威利，1972年。

[Bezi86] Bézier, P.E.，《UNISURF CAD系统的数学基础》，伦敦:巴特沃斯，1986年。

[Blom91] Bloomenthal, M. 和 Riesenfeld, R.F.，通过张量积NURBS近似扫掠曲面，计算机视觉与图形中的曲线与表面 II，SPIE 会议论文集(光学仪器工程师学会)，第1610卷，页码132-144，1991年。

[Blor91] Bloor, M. 和 Owen, J.，CAD/CAM产品数据交换:下一步，CAD，第23卷，第4期，页码237-243，1991年。

[Boeh80] Boehm, W.，在B样条曲线中插入新节点，CAD，第12卷，第4期，页码199-201，1980年。

[Boeh84] Boehm, W.，Farin, G. 和 Kahmann, J.，CAGD中的曲线和表面方法综述，计算机辅助几何设计，第1卷，第1期，页码1-60，1984年。

[Boeh85a] Boehm, W. 和 Prautzsch, H.，插入算法，CAD，第17卷，第2期，页码58-59，1985年。

[Boeh85b] Boehm, W., 关于结点插入算法的效率, 计算机辅助几何设计, 第2卷，第1-3期，页码141-143, 1985年.

[Boeh91] Boehm, W., 和 Hansford, D., 二次曲面上的Bézier片段, 收录于《曲线和曲面设计中的NURBS》，Farin, G., 编辑，费城: SIAM, 页码1-14, 1991年.

[Boeh94] Boehm, W., 和 Prautzsch, H., 几何设计的几何概念, Wellesley, MA: A.K. Peters, 1994年.

[Bron92] Bronsvoort, W., 和 Waarts, J., 一种将广义圆柱的表面转换为B样条表面的方法, 计算与图形, 第16卷，第2期，页码175-178, 1992年.

[Butt76] Butterfield, K.R., B样条基的所有导数的计算, 数学应用研究所期刊, 第17卷，页码15-25, 1976年.

[Chan81] Chang, G., 和 Wu, J., Bézier技术的数学基础, CAD, 第13卷，第3期，页码133-136, 1981年.

[Choi90] Choi, B.K., 和 Lee, C., 通过坐标变换和混合进行扫掠曲面建模, CAD, 第22卷，第2期，页码87-96, 1990年.

[Choi91] Choi, B.K., CAD/CAM的表面建模, 纽约: Elsevier, 1991年.

[Chou92] Chou, J., 和 Piegl, L., 使用三次有理 B-spline 进行数据缩减, IEEE 计算机图形与应用, 第 12 卷, 第 3 期, 第 60-68 页, 1992 年.

[Chou95] Chou, J.J., 高阶 Bézier 圆, CAD, 即将出版, 1995 年.

[Cobb84] Cobb, E.S., "使用 B-spline 表示法设计雕塑表面," 博士学位论文, 犹他大学, 1984 年.

[Cohe80] Cohen, E., Lyche, T., 和 Riesenfeld, R.F., 离散 B-spline 和细分技术在计算机辅助几何设计和计算机图形学中的应用, 计算机图形与图像处理, 第 14 卷, 第 87-111 页, 1980 年.

[Cohe85] Cohen, E., Lyche, T., 和 Schumaker, L.L., 提高样条曲线次数的算法, ACM TOG, 第 4 卷, 第 3 期, 第 171-181 页, 1985 年.

[Coon67] Coons, S.A., 用于计算机辅助设计空间形状的表面, MAC-TR-41, 麻省理工学院, 1967 年 6 月.

[Coqu87a] Coquillart, S., 一种基于控制点的扫掠技术, IEEE 计算机图形与应用, 第 7 卷, 第 11 期, 第 36-45 页, 1987 年.

[Coqu87b] Coquillart, S., 计算 B-spline 曲线的偏移, CAD, 第 19 卷, 第 6 期, 第 305-309 页, 1987 年.

[Coqu90] Coquillart, S., 扩展自由形变:一种用于三维几何建模的雕刻工具, Comput. Graph., 第24卷, 第4期, 页187-193, 1990年.

[Coqu91] Coquillart, S., 和 Jancene, P., 动画自由形变:一种交互式动画技术, Comput. Graph., 第25卷, 第4期, 页23-26, 1991年.

[Cox72] Cox, M.G., B样条的数值评估, Jour. Inst. Math. Ap-plic., 第10卷, 页134-149, 1972年.

[Coxe67] Coxeter, H.S.M., 和 Greitzer, S.L., 几何学再探, 华盛顿特区:美国数学协会, 1967年.

[Coxe74] Coxeter, H.S.M., 投影几何, 加拿大多伦多:多伦多大学出版社, 1974年.

[Coxe80] Coxeter, H.S.M., 几何学导论, 纽约:约翰·威利出版社, 1980年.

[Curr47] Curry, H.B., 和 Schoenberg, I.J., 关于样条分布及其极限:波利亚分布函数, 摘要380t, 美国数学学会公报, 第53卷, 页109, 1947年.

[Dani89] Daniel, M., 和 Daubisse, J.C., 在幂基中使用贝塞尔曲线和曲面的数值问题, 计算辅助几何设计, 第6卷, 页121-128, 1989年.

[Dann85] Dannenberg, L., 和 Nowacki, H., 使用多项式基进行表面表示的近似转换, 计算机辅助几何设计, 第2卷, 第123-132页, 1985年.

[DeBo72] De Boor, C., 关于使用 B-spline 进行计算, 近似理论杂志, 第6卷, 第50-62页, 1972年.

[DeBo78] De Boor, C., B-spline 实用指南, 纽约: Springer-Verlag, 1978年.

[DeBo87] De Boor, C., 切角总是有效, 计算机辅助几何设计, 第4卷, 第1-2期, 第125-131页, 1987年.

[DeBo93] De Boor, C., B(asic)-spline 基础, 收录于 [Pieg93].

[deCa86] de Casteljau, P., 形状数学与计算机辅助设计, 伦敦: Kogan Page, 1986年.

[deCa93] de Casteljau, P., 用于曲线和曲面建模的极坐标形式，应用于雪铁龙, 收录于 [Pieg93].

[deMo84] de Montaudouin, Y., 和 Tiller, W., 计算机辅助几何设计中的 Cayley 方法, 计算机辅助几何设计, 第1卷, 第4期, 第309-326页, 1984年.

[DoCa76] Do Carmo, M.P., 曲线和曲面的微分几何, Englewood Cliffs, NJ: Prentice-Hall, 1976年.

[Eck93] Eck, M., Bézier 曲线的度数降低, 计算机辅助几何设计, 第10卷, 第237-251页, 1993年.

[Fari83] Farin, G.E., 有理 Bézier 曲线的算法, CAD, 第 15 卷, 第 2 期, 第 73-77 页, 1983 年.

[Fari89] Farin, G.E., 有理曲线和曲面, 收录于《计算机辅助几何设计中的数学方面》，Lyche, T. 和 Schumaker, L.L. 主编, 纽约: 学术出版社, 1989 年.

[Fari93] Farin, G.E., 计算机辅助几何设计的曲线与曲面 - 实用指南, 第 3 版, 波士顿: 学术出版社, 1993 年.

[Faro86] Farouki, R., 非退化偏移曲面的逼近, 计算机辅助几何设计, 第 3 卷, 第 15-43 页, 1986 年.

[Faro87] Farouki, R.T., 和 Rajan, V.T., 关于伯恩斯坦形式多项式的数值条件, 计算机辅助几何设计, 第 4 卷, 第 191-216 页, 1987 年.

[Faro88] Farouki, R.T., 和 Rajan, V.T., 伯恩斯坦形式多项式的算法, 计算机辅助几何设计, 第 5 卷, 第 1-26 页, 1988 年.

[Faux81] Faux, I.D., 和 Pratt, M.J., 设计与制造的计算几何, 英国奇切斯特: Ellis Horwood Ltd., 1981 年.

[Ferg66] Ferguson, J.C., 在特定参数曲线类中表征的形状, 报告 3122-31, TRW 公司, 加利福尼亚州雷东多海滩, 1966 年.

[Ferg67] Ferguson, J.C., 以特殊类别的参数曲线为特征的形式 - II，报告 3122-3-237，TRW 公司，加州雷东多海滩，1967年。

[Ferg69] Ferguson, J.C. 和 Miller, K.L., 在一类三次代数曲线中对形状的表征，报告 5322-3-5，TRW 公司，加州雷东多海滩，1969年。

[Ferg93] Ferguson, J.C., 自由形状曲线和超曲面的 F 方法定义，见 [Pieg93]。

[Fili89] Filip, D. 和 Ball, T., 程序化表示抬升表面，IEEE 计算机图形与应用，第 9 卷，第 6 期，页 27-33，1989年。

[Fole90] Foley, J., van Dam, A., Feiner, S. 和 Hughes, J., 计算机图形学:原理与实践，马萨诸塞州雷丁:阿迪森-韦斯利，1990年。

[Forr68] Forrest, A.R., “计算机辅助设计的曲线与表面，”博士论文，剑桥大学，英国剑桥，1968年。

[Forr70] Forrest, A.R., 非理性扭曲立方曲线的形状分类，基于 Bézier 多边形，CAD 组文件编号 52，剑桥大学，英国剑桥，1970年。

[Forr72] Forrest, A.R., 通过 Bézier 多项式进行交互式插值和逼近，计算机期刊，第 15 卷，第 1 期，页 71-79，1972年。更正和更新版本见 \({CAD}\) ，第 22 卷，第 9 期，页 527-537，1990年。

[Forr80] Forrest, A.R., 扭曲立方曲线:一种计算机辅助几何设计方法, CAD, 第12卷，第2期，页165-172, 1980.

[Fors88] Forsey, D., 和 Bartels, R., 层次B样条细化, Comput. Graph., 第22卷，第4期，页205-212, 1988.

[Four83] Fournier, A., 和 Wesley, M., 弯曲多面体对象, CAD, 第15卷，第2期，页79-87, 1983.

[Fowl92] Fowler, B., 张量积曲面的几何操作, Comput. Graph. 特刊, 交互式3D图形研讨会, 页101-108, 1992.

[Fowl93] Fowler, B., 和 Bartels, R., 基于约束的曲线操作, IEEE Comput. Graph. and Appl., 第13卷，第5期，页43-49, 1993.

[Fuhr92] Fuhr, R.D., 和 Kallay, M., 单调线性有理样条插值, Comput. Aid. Geom. Des., 第9卷，第4期，页313-319, 1992.

[Geis90] Geise, G., 和 Langbecker, U., 具有四条圆锥边界曲线的有限二次段, Comput. Aid. Geom. Des., 第7卷，页141-150, 1990.

[Good81] Goodman, T., 通过Hermite-Birkhoff样条进行Hermite-Birkhoff插值, Proc. Roy. Soc. Edinburgh, 第88卷(A)，第3/4部分，页195-201, 1981.

[Gord69] Gordon, W., 通过曲线网络进行样条混合表面插值, 数学与力学杂志, 第18卷, 第10期, 第931-952页, 1969年.

[Gord71] Gordon, W., 二元和多元插值与逼近的混合函数方法, SIAM 数值分析杂志, 第8卷, 第158-177页, 1971年.

[Gord74a] Gordon, W.J., 和 Riesenfeld, R.F., 用于计算机辅助设计自由形状曲线和表面的伯恩斯坦-贝塞尔方法, 计算机协会杂志, 第21卷, 第2期, 第293-310页, 1974年.

[Gord74b] Gordon, W.J., 和 Riesenfeld, R.F., B样条曲线和表面, 收录于《计算机辅助几何设计》，Barnhill, R.E., 和 Riesenfeld, R.F., 编辑, 纽约: 学术出版社, 1974年.

[Gord93] Gordon, W., 通过混合函数方法定义雕刻表面, 收录于 [Pieg93].

[Gugg89] Guggenheimer, H., 沿轨迹计算框架, 计算机辅助几何设计, 第6卷, 第77-78页, 1989年.

[Hoff89] Hoffmann, C.M., 几何与实体建模, 加州圣马特奥: 摩根考夫曼, 1989年.

[Hohm91] Hohmeyer, M. 和 Barsky, B., 皮肤化有理B样条曲线以构建插值表面, 计算机视觉、图形与图像处理: 图形模型与图像处理, 第53卷, 第6期, 第511-521页, 1991年.

[Hosc88] Hoschek, J., 偏移曲线的样条逼近, Comput. Aid. Geom. Des., 第 5 卷, 第 33-40 页, 1988.

[Hosc92a] Hoschek, J., 二次曲面上的贝塞尔曲线和曲面片, 收录于《计算机辅助几何设计中的数学方法 II》，Lyche, T. 和 Schumaker, L. 编辑, 纽约: 学术出版社, 第 331-342 页, 1992.

[Hosc92b] Hoschek, J., 圆形样条, CAD, 第 24 卷, 第 11 期, 第 611-618 页, 1992.

[Hosc93] Hoschek, J. 和 Lasser, D., 计算机辅助几何设计基础, Wellesley, MA: A.K. Peters, Ltd., 1993.

[Howa91a] Howard, T., 评估 PHIGS 在 CAD 和一般图形应用中的表现, CAD, 第 23 卷, 第 4 期, 第 244-251 页, 1991.

[Howa91b] Howard, T.L.J., Hewitt, W.T., Hubbold, R.J., 和 Wyrwas, K.M., PHIGS 和 PHIGS PLUS 的实用介绍, Reading, MA: Addison-Wesley, 1991.

[Hsu92] Hsu, W., Hughes, J., 和 Kaufman, H., 自由形状变形的直接操作, Comput. Graph., 第 26 卷, 第 2 期, 第 177-184 页, 1992.

[IGE93] 初始图形交换规范 (IGES) 第 5.2 版, ANSI Y14.26M, 可从美国产品数据协会 (US PRO) 获取, Fairfax, VA, USA, 1993.

[Ilyi84] Ilyin, V.A., 和 Poznyak, E.G., 《解析几何》，莫斯科:Mir Publishers，1984年。

[Kapl52] Kaplan, W., 《高级微积分》，马萨诸塞州雷丁:Addison Wesley，1952年。

[Klok86] Klok, F., 两个移动坐标框架用于沿3D轨迹的扫掠，计算机辅助几何设计，卷3，第217-229页，1986年。

[Knut73] Knuth, D.E., 《计算机程序设计的艺术》。第1卷，基础算法，马萨诸塞州雷丁:Addison-Wesley，1973年。

[Lach88] Lachance, M.A., 参数曲面的切比雪夫经济化，计算机辅助几何设计，卷5，第195-208页，1988年。

[Lanc86] Lancaster, P., 和 Salkauskas, K., 《曲线和曲面拟合》，纽约:学术出版社，1986年。

[Lane80] Lane, J.M., 和 Riesenfeld, R.F., 一种用于计算机生成和显示分段多项式曲面的理论发展，IEEE模式分析与机器智能杂志，卷PAMI-2，第1期，第35-46页，1980年。

[Lane83] Lane, J.M., 和 Riesenfeld, R.F., B样条逼近的变差减小性质的几何证明，逼近理论杂志，卷37，第1-4页，1983年。

[Laur93] Laurent-Gengoux, P., 和 Mekhilef, M., NURBS表示的优化，计算机辅助设计，卷25，第11期，第699-710页，1993年。

[Lawr72] Lawrence, J.D., 特殊平面曲线目录, 纽约: Dover, 1972.

[Lee83] Lee, E.T.Y., B-spline 入门, 波音文档, 1983.

[Lee87] Lee, E.T.Y., 用于圆锥曲线的有理二次Bézier表示, 收录于《几何建模:算法与新趋势》，Farin, G.E. 编辑，费城: SIAM, 第3-19页, 1987.

[Lee89] Lee, E.T.Y., 在参数曲线插值中选择节点, CAD, 第21卷, 第363-370页, 1989.

[Lee91] Lee, E.T.Y. 和 Lucian, M.L., 有理B样条的Möbius重新参数化, 计算机辅助几何设计, 第8卷, 第213-215页, 1991.

[Limi44] Liming, R.A., 实用解析几何及其在航空器中的应用, 纽约: Macmillan, 1944.

[Limi79] Liming, R.A., 计算机图形学数学, Fallbrook, CA: Aero Publishers Inc., 1979.

[Lin94] Lin, F. 和 Hewitt, W., 将Coons-Gordon曲面表示为NURBS, CAD, 第26卷, 第2期, 第145-155页, 1994.

[Lore86] Lorentz, G.G., Berntein多项式, 纽约: Chelsea Publishing Co., 1986.

[Lych85] Lyche, T., Cohen, E. 和 Morken, K., 张量积样条的节点线细化算法, 计算机辅助几何设计, 第2卷, 第1-3期, 第133-139页, 1985.

[Lych87] Lyche, T., 和 Morken, K., 参数 B-spline 曲线和曲面的节点去除, Comput. Aid. Geom. Des., 第 4 卷, 第 217-230 页, 1987.

[Lych88] Lyche, T., 和 Morken, K., 一种用于样条的数据减少策略及其在函数和数据近似中的应用, IMA Jour. Num. Anal., 第 8 卷, 第 185-208 页, 1988.

[Mats92] Matsuki, N., 一种用于 B-spline 曲面的交互式形状修改方法, 收录于计算机集成制造中的人类因素, Olling, G., 和 Kimura, F., 编辑, 阿姆斯特丹: Elsevier Science Publishers B.V. (North-Holland) (IFIP), 第 385-397 页, 1992.

[Mort85] Mortenson, M.E., 几何建模, 纽约: John Wiley, 1985.

[Patt85] Patterson, R.R., 伯恩斯坦-贝济耶曲线参数的投影变换, ACM TOG, 第 4 卷, 第 4 期, 第 276-290 页, 1985.

[Pieg86] Piegl, L., 计算机辅助设计中有理贝济耶方案的几何研究, Comput. in Industry, 第 7 卷, 第 401-410 页, 1986.

[Pieg87a] Piegl, L., 和 Tiller, W., 使用有理 B-splines 的曲线和曲面构造, CAD, 第 19 卷, 第 9 期, 第 485-498 页, 1987.

[Pieg87b] Piegl, L., 一种使用圆锥曲线平滑散乱数据的技术, Comput. in Industry, 第 9 卷, 第 223-237 页, 1987.

[Pieg87c] Piegl, L., 关于在计算机辅助几何设计中使用无限控制点, Comput. Aid. Geom. Des., 第 4 卷, 第 155-166 页, 1987 年.

[Pieg87d] Piegl, L., 通过有理 Bézier 曲线进行交互式数据插值, IEEE Comput. Graph. and Appl., 第 7 卷, 第 4 期, 第 45-58 页, 1987 年.

[Pieg88a] Piegl, L., 使用具有无限控制点的有理 Bézier 逼近形式的 Hermite 和 Coons 类插值, CAD, 第 20 卷, 第 1 期, 第 2-10 页, 1988 年.

[Pieg88b] Piegl, L., Coons 类型的补丁, Comput. and Graph., 第 12 卷, 第 2 期, 第 221-228 页, 1988 年.

[Pieg89a] Piegl, L., 计算机辅助几何设计中的关键发展, CAD, 第 21 卷, 第 5 期, 第 262-273 页, 1989 年.

[Pieg89b] Piegl, L., 和 Tiller, W., 有理 B 样条圆的动物园, IEEE Comput. Graph. and Appl., 第 9 卷, 第 5 期, 第 48-56 页, 1989 年.

[Pieg89c] Piegl, L., 修改有理 B 样条的形状. 第 1 部分: 曲线, CAD, 第 21 卷, 第 8 期, 第 509-518 页, 1989 年.

[Pieg89d] Piegl, L., 修改有理 B 样条的形状. 第 2 部分: 表面, CAD, 第 21 卷, 第 9 期, 第 538-546 页, 1989 年.

[Pieg90] Piegl, L., 圆锥样条计算算法, 土木工程计算杂志, 第4卷, 第3期, 页180-197, 1990年.

[Pieg91a] Piegl, L., 关于NURBS的调查, IEEE计算机图形与应用, 第10卷, 第1期, 页55-71, 1991年.

[Pieg91b] Piegl, L., 和 Tiller, W., B样条曲线的存储高效分解, CSE 91-01, 南佛罗里达大学计算机科学与工程系, 坦帕, FL, 1991年.

[Pieg93] Piegl, L., 编辑, 计算机辅助几何建模的基础发展, 伦敦: 学术出版社, 1993年.

[Pieg94] Piegl, L., 和 Tiller, W., B样条曲线的度提升的软件工程方法, CAD, 第26卷, 第1期, 页17-28, 1994年.

[Pieg95] Piegl, L., 和 Tiller, W., B样条曲线的度降低算法, CAD, 第27卷, 第2期, 1995年.

[Prau84] Prautzsch, H., B样条曲线的度提升, 计算机辅助几何设计, 第1卷, 第1期, 页193-198, 1984年.

[Prau91] Prautzsch, H., 和 Piper, B., 提高样条曲线度的快速算法, 计算机辅助几何设计, 第8卷, 页253-265, 1991年.

[Prau92] Prautzsch, H., 和 Gallagher, T., B样条或贝塞尔曲面是否具有几何变异减小特性？，计算机辅助几何设计，卷 9，第 2 期，页 119-124，1992年。

[Pres88] Press, W., Flannery, B., Teukolsky, S., 和 Vetterling, W., 数值计算方法 \(C\) ，剑桥，英国:剑桥大学出版社，1988年。

[PHIG92] 程序员的分层交互图形系统 (PHIGS)，ISO/IEC 9592-4: (1992)，可从美国国家标准与技术研究院 (NIST) 获取，马里兰州盖瑟斯堡。

[Rams87] Ramshaw, L., 绽放:一种连接点的样条方法，报告 19，数字系统研究中心，加州帕洛阿尔托，1987年。

[Renn82] Renner, G., 机械工程实践中的形状描述方法，工业计算，卷 3，页 137-142，1982年。

[Ries73] Riesenfeld, R.F., “B样条逼近在计算机辅助设计几何问题中的应用，”博士论文，锡拉丘兹大学，1973年。

[Ries81] Riesenfeld, R.F., 计算机图形学中的齐次坐标和射影平面，IEEE 计算机图形与应用，卷 1，第 1 期，页 50-55，1981年。

[Robe65] Roberts, L.G., 齐次矩阵表示和 \(n\) 维构造的操作，技术报告 MS-1405，林肯实验室，麻省理工学院，马萨诸塞州剑桥，1965年。

[Roge90] Rogers, D.F., 和 Adams, J.A., 计算机图形学的数学元素, 第2版, 纽约: McGraw-Hill, 1990.

[Salm79] Salmon, G., 圆锥曲线论, Longman, Green \& Co., 第6版, 伦敦, 1879, 由 Dover 出版社重印, 纽约.

[Scho46] Schoenberg, I.J., 关于用解析函数逼近等距数据的问题的贡献, Quart. Appl. Math., 第4卷, 第45-99页, 1946.

[Scho66] Schoenberg, I.J., 关于 Hermite-Birkhoff 插值, 数学分析与应用杂志, 第16卷, 第3期, 第538-543页, 1966.

[Sede84] Sederberg, T., Anderson, D., 和 Goldman, R., 参数曲线和曲面的隐式表示, 计算机视觉、图形与图像处理, 第28卷, 第72-84页, 1984.

[Sede86] Sederberg, T., 和 Parry, S., 实体几何对象的自由形变, 计算机图形, 第20卷, 第4期, 第151-160页, 1986.

[Silt92] Siltanen, P., 和 Woodward, C., 3D 偏移曲线、扫掠曲面和蒙皮的法向方向方法, 欧洲图形学会议论文集 92, 第11卷, 第3期, 第C-449-C-457页, 1992.

[Smit71] Smith, L.B., 用固定数量的点和最大内切面积绘制椭圆、双曲线或抛物线, 计算机期刊, 第14卷，第1期，81-86页，1971年。

[Smit74] Smith, R., Price, J., 和 Howser, L., 使用三次样条函数的平滑算法, NASA技术报告, TN D-7397, NASA兰利研究中心, 亨普顿, VA, 1974年。

[STEP94] 产品模型数据交换标准 (STEP), ISO 10303, 一系列文档。第42部分:几何和拓扑表示, ISO 10303-42, 可从ISO秘书处获取, 国家标准与技术研究所 (NIST), 盖瑟斯堡, MD, 1994年。

[Ston89] Stone, M.C., 和 DeRose, T.C., 参数立方曲线的几何特征, ACM TOG, 第8卷，第3期，147-163页，1989年。

[Stra92] Straker, D., C风格标准与指南, 纽约:普伦蒂斯霍尔, 1992年。

[Stro91] Strotman, T., 私人通信, 1991年。638 参考文献

[Su89] Su, B., 和 Liu, D., 计算几何 - 曲线与曲面建模, 波士顿:学术出版社, 1989年。

[Till83] Tiller, W., 用于曲线和曲面表示的有理B样条, IEEE计算机图形与应用, 第3卷，第6期，61-69页，1983年。

[Till84] Tiller, W. 和 Hanson, E., 二维轮廓的偏移，IEEE 计算机图形与应用，卷 4，第 9 期，页 36-46，1984 年。

[Till92] Tiller, W., NURBS 曲线和曲面的去结算法，CAD，卷 24，第 8 期，页 445-453，1992 年。

[Vand83] Vandergraft, J., 数值计算导论，纽约:学术出版社，1983 年。

[Verg91] Vergeest, J., 使用 STEP 的 CAD 表面数据交换，CAD，卷 23，第 4 期，页 269-281，1991 年。

[Vers75] Versprille, K.J., “有理 B 样条逼近形式的计算机辅助设计应用，”博士论文，锡拉丘兹大学，1975 年。

[Wang81] Wang, C.Y., 参数立方曲线和参数 B 样条立方曲线的形状分类，CAD，卷 13，第 4 期，页 199-206，1981 年。

[Watk88] Watkins, M.A. 和 Worsey, A.J., Bézier 曲线的度数降低，CAD，卷 20，第 7 期，页 398-405，1988 年。

[Wein92] Weinstein, S.E. 和 Xu, Y., 通过逼近和插值降低 Bézier 曲线的度数，见《逼近理论》，Anastassiou, G. A. 编辑，纽约:Dekker，1992 年，页 503-512。

[Welc92] Welch, W. 和 Witkin, A.，变分表面建模，计算机图形，卷 26，第 2 期，页 157-166，1992 年。

[Wood87] Woodward, C.，B-spline 表面的横截面设计，计算机与图形，卷 11，第 2 期，页 193-201，1987 年。

[Wood88] Woodward, C.，用于交互式 B-spline 表面插值的皮肤技术，计算机辅助设计，卷 20，第 8 期，页 441-451，1988 年。

[Yama88] Yamaguchi, F.，计算机辅助几何设计中的曲线和表面，纽约:施普林格出版社，1988 年。



\clearpage
    

\end{document}


